% =====================================================================
\chapter{Kiértékelési módszertan}
\label{ch:kiertekeles}
% =====================================================================

\Az{\ref{ch:szakirodalom}}.~fejezetben láttuk, hogy a kérdésgenerálás kiértékelése jellemzően a
referencia-átfedésre és kis, nem vak humán értékelésekre támaszkodik, és ritkán ellenőrzi a kulcs
helyességét. Ebben a fejezetben olyan kiértékelési módszertant mutatunk be, amely (i) \az{\ref{ch:formalizmus}}.~fejezetben definiált, forrásalapú predikátumokat méri, (ii) egytényezős
kísérleti tervvel választja szét a komponensek hatását, (iii) a mérőeszközt (a bírót) maga is
validálja, és (iv) a dokumentumot tekinti a statisztikai elemzés egységének. A módszertan
minden eleme a nyílt forráskódú keretrendszerben (\code{tests/eval}) implementált.

\section{A kiértékelő keretrendszer}
\label{sec:harness}

A keretrendszer egy parancssori eszköz (\code{python -m mimir\_eval}), amelynek fő lépései \az{\ref{fig:harness}}.~ábrán láthatók. A \code{run} lépés egy kar konfigurációja szerint minden
dokumentum\,$\times$\,véletlenmag párra vizsgát generál, és a vizsgákat a teljes nyers kimenettel,
a visszakeresett kontextussal, a hívásstatisztikával és a GPU-memória-mintákkal együtt JSONL-fájlba
írja. A \code{score} lépés kiszámítja az automatikus mutatókat és meghívja a bírót. A \code{report} lépés
táblázatokat, ábrákat, szignifikanciateszteket és a dolgozat \LaTeX{}-táblázatait állítja elő, az
\code{analyze} lépés pedig a karpárok összehasonlítását végzi. A \code{rate-export} és \code{rate-import}
lépések az oktatói értékelőlapokat kezelik.

\begin{figure}[htbp]
\centering
\resizebox{\textwidth}{!}{%
\begin{tikzpicture}[
  node distance=5mm and 5mm,
  box/.style={draw, rounded corners=2pt, align=center, font=\footnotesize, minimum height=12mm, text width=27mm, fill=blue!6, draw=blue!50!black},
  file/.style={draw, align=center, font=\scriptsize, minimum height=9mm, text width=31mm, fill=gray!8, rounded corners=1pt},
  arr/.style={-{Stealth[length=1.6mm]}, semithick}]
\node[file] (cfg) {kar konfigurációja\\(YAML) és\\\texttt{manifest.yaml}};
\node[box, right=of cfg] (run) {\texttt{run}\\generálás:\\dokumentum $\times$ mag};
\node[box, right=of run] (score) {\texttt{score}\\automatikus mutatók\\és LLM-bíró};
\node[box, right=of score] (report) {\texttt{report}\\táblák, ábrák,\\tesztek, \LaTeX};
\node[file, below=of run] (ex) {\texttt{exams.jsonl}\\\texttt{run.json}};
\node[file, below=of score] (q) {\texttt{questions.jsonl}\\\texttt{exam\_metrics.csv}\\\texttt{judge\_cache.jsonl}};
\node[file, below=of report] (tex) {\texttt{tab\_*.tex}\\\texttt{paper\_numbers.json}};
\node[box, above=of score, fill=orange!8, draw=orange!70!black] (rate) {\texttt{\small rate-export/}\\\texttt{\small import}: oktatói értékelés};
\draw[arr] (cfg) -- (run); \draw[arr] (run) -- (score); \draw[arr] (score) -- (report);
\draw[arr] (run) -- (ex); \draw[arr] (score) -- (q); \draw[arr] (report) -- (tex);
\draw[arr] (score) -- (rate); \draw[arr] (rate) -| (report);
\end{tikzpicture}}
\caption{A kiértékelő keretrendszer lépései és kimenetei.}
\label{fig:harness}
\end{figure}

A keretrendszer a kinyerést, a darabolást és a visszakeresést a futó szolgáltatások HTTP-végpontjain
keresztül végzi, így pontosan azt méri, ami éles üzemben fut; a naiv alapvonal promptját a Bifrost
függőségmentes \code{prompts} moduljából importálja, így az bájtra azonos az éles prompttal; a tervezett
csővezetéket pedig konfigurációs kulcsokkal kapcsolható komponensekkel valósítja meg. Minden hívás mellé
rögzíti a hívás szakaszát, a tokenszámokat és a késleltetést, a GPU-memóriát pedig egy háttérszál
mintavételezi. A reprodukálhatóságot négy mechanizmus biztosítja. (1) Minden futás rögzíti a feloldott
konfigurációt, a git-verziót (a „piszkos” munkakönyvtár jelzésével), az adatkészlet verzióját, a
generátor beállításait és a hardvert (GPU neve, teljes memóriája, meghajtóverzió). (2) A bíró
válaszai a (bírómodell, promptverzió, gondolkodási szint, prompt) négyes kriptográfiai lenyomata szerint
gyorsítótárba kerülnek, így az újrapontozás determinisztikus, és csak az új kérdésekért kell fizetni.
(3) A futások megszakítás után folytathatók: a már elkészült dokumentum\,$\times$\,mag párokat a
keretrendszer kihagyja. (4) A teljes \code{run}$\to$\code{score}$\to$\code{report} folyamatot hamis
szolgáltatásokkal és ál-nyelvi modellel futó végponttól végpontig tartó tesztek fedik le
(a keretrendszer egységtesztjei, 40 tesztfüggvény).

\section{Adatkészlet}

A kiértékelő adatkészlet (\code{eval-v1}) 39 dokumentumból áll, három forrásból
(\ref{tab:adatkeszlet}.~táblázat). A forrásokat az döntötte el, hogy minden kérdéshez rendelkezésre
álljon a forrásszakasz \emph{és} a választ igazoló mondat (\emph{evidence}), mert a visszakeresési és
tematikus mutatók ezt igénylik.

\begin{table}[htbp]
\centering
\caption{Az \code{eval-v1} kiértékelő adatkészlet. R/K/H: rövid ($\le$\,6\,000), közepes
(6\,000--25\,000) és hosszú (25\,000--90\,000 karakter) dokumentumok száma.}
\label{tab:adatkeszlet}
\small
\begin{tabular}{@{}llrcrrl@{}}
\toprule
Forrás & Nyelv & Dok. & R/K/H & Medián kar. & Ref.\ kérdés & Licenc \\
\midrule
Saját oktatási szövegek & HU & 3 & 3/0/0 & 1\,923 & 12 & saját \\
OpenStax, EduQG~\cite{hadifar2023eduqg} & EN & 12 & 4/4/4 & 16\,752 & 74 & CC BY 4.0 \\
Wikipédia, MILQA~\cite{novak2023milqa} & HU & 24 & 8/8/8 & 8\,072 & 240 & CC BY-SA 4.0 \\
\midrule
Összesen & & 39 & 15/12/12 & 7\,919 & 326 & \\
\bottomrule
\end{tabular}
\end{table}

\paragraph{Saját szövegek.} Három rövid magyar oktatási szöveg (a kávé története és feldolgozása;
az immunrendszer; a játékelmélet alapjai), könnyű, közepes és nehéz vizsgaszinttel, dokumentumonként négy,
kézzel írt referenciakérdéssel és kézzel megjelölt bizonyító mondatokkal. A pilot vizsgálat ezek
közül az elsőn (\code{hu-coffee}) futott.

\paragraph{EduQG.} Az EduQG adatkészlet~\cite{hadifar2023eduqg} OpenStax-tankönyvekhez írt, szakértői
feleletválasztós kérdéseket tartalmaz bizonyító mondatokkal és részben Bloom-címkékkel. Az importáló
tárgyterületenként (élettudományok, társadalomtudományok, jog, közgazdaságtan) egy rövid, egy közepes és
egy hosszú dokumentumot választ: összefüggő bekezdésablakot keres, amely a hossztartományba esik és a
lehető legtöbb referenciakérdés bizonyító mondatát tartalmazza; csak azokat a kérdéseket tartja meg,
amelyek bizonyító mondata szó szerint szerepel a dokumentumban.

\paragraph{MILQA.} A MILQA~\cite{novak2023milqa} a legnagyobb magyar szövegértési adatbázis rövid és
hosszú válaszszakaszokkal. A hosszú dokumentumok a Wikipédia azonos című, aktuális cikkéből
készülnek; referenciakérdésként azokat tartjuk meg, amelyek bizonyító mondata a jelenlegi cikkben is
megtalálható.

A referenciakérdéseket \emph{nem} használjuk célként: egy jó vizsgakérdésnek nem kell a
referenciakérdéshez hasonlítania. Szerepük kettős: bizonyító mondataik alapján mérjük a visszakeresés
pontosságát (recall@$k$, MRR), valamint azt, hogy a generált kérdések a referencia által fontosnak tartott
tényeket érintik-e. Az adatkészletet az első éles futás előtt \code{eval-v1} verzióként befagyasztjuk;
minden későbbi változtatás új verziót kap, és a futások rögzítik a verziót, így eltérő verziókon
készült eredmények nem keveredhetnek.

\section{Kísérleti karok}
\label{sec:karok}

A kísérleti terv \emph{egytényezős}: minden kar egy referenciakarhoz képest pontosan egy tényezőben
tér el, így a különbség az adott tényezőnek tulajdonítható. \Az{\ref{tab:karok}}.~táblázat a karokat, \az{\ref{tab:osszevetesek}}.~táblázat a tervezett páros összehasonlításokat és azok kutatási kérdéseit
sorolja fel. Minden helyi kar a Qwen2.5-7B-Instruct modellt használja Q4\_K\_M kvantálással, $T=0{,}2$
hőmérséklettel; a szerverkarok az egyetemi GenAI-szerver nagy Qwen-családba tartozó modelljét
(\code{qwen38udq8xl}).

\begin{table}[htbp]
\centering
\caption{Kísérleti karok. Minden kar a zárójelben megadott referenciakartól egy tényezőben tér el.}
\label{tab:karok}
\small
\begin{tabularx}{\textwidth}{@{}l>{\raggedright\arraybackslash}Xc@{}}
\toprule
Kar & Konfiguráció & Magok \\
\midrule
\arm{B-doc-L} & teljes dokumentum egy promptban, helyi 7B, 16k kontextus (status quo, helyi) & 3 \\
\arm{B-doc-S} & teljes dokumentum egy promptban, szervermodell (status quo) & 3 \\
\arm{E0}  & naiv RAG: 3 darab, a teljes vizsga egy hívásban (alapvonal) & 3 \\
\arm{E0s} & mint \arm{E0}, de a Bifrost éles végpontján keresztül (implementációs kontroll) & 3 \\
\arm{E1}  & vizsgaterv: tervező, slotonkénti keresés, kérdésenként egy hívás (\arm{E0}) & 1 \\
\arm{E2}  & \arm{E1} + ellenőrző lánc visszacsatolással (\arm{E1}) & 1 \\
\arm{E2f} & \arm{E2} vak megoldási próba nélkül (\arm{E2}) & 1 \\
\arm{E2h} & \arm{E2} hibrid (BM25 + sűrű, RRF) visszakereséssel (\arm{E2}) & 1 \\
\arm{E2x} & \arm{E2} más modellcsaládba tartozó ellenőrzővel, Llama-3.1-8B (\arm{E2}) & 1 \\
\arm{E3}  & \arm{E2} + fogalomgráf (\arm{E2}) & 1 \\
\arm{E4}  & \arm{E0} a szervermodellel (\arm{E0}) & 3 \\
\arm{E4b} & \arm{E2} a szervermodellel (\arm{E2}) & 1 \\
\arm{E5a} & \arm{E0} a régi, 512 tokenre csonkolt kódolóval (\arm{E0}) & 3 \\
\arm{E5b} & \arm{E0} $\mu+1{,}5\sigma$ küszöbbel (\arm{E0}) & 3 \\
\arm{E5c} & \arm{E0} rögzített, 800 karakteres darabokkal (\arm{E0}) & 3 \\
\bottomrule
\end{tabularx}
\end{table}

A tervezett karok (\arm{E1--E3}, \arm{E4b}) egy véletlenmaggal futnak, mert vizsgánként 3--8 percet
igényelnek, és a variancia elsődleges forrása a dokumentum, nem a mag; az egyhívásos karok három
maggal futnak, mert olcsók, és a pilot szerint (\ref{sec:zaj}.~szakasz) egyetlen futás zaja jelentős.

\begin{table}[htbp]
\centering
\caption{Tervezett páros összehasonlítások és a hozzájuk tartozó kutatási kérdések, hipotézisek.}
\label{tab:osszevetesek}
\small
\begin{tabular}{@{}llll@{}}
\toprule
Összevetés & Vizsgált tényező & RQ & Hipotézis \\
\midrule
\arm{E1} -- \arm{E0} & tervezés & RQ1, RQ2 & \ref{h:terv} \\
\arm{E2} -- \arm{E1} & ellenőrzés & RQ2 & \ref{h:verif} \\
\arm{E2x} -- \arm{E2} & ellenőrző modellcsaládja & RQ2 & \ref{h:verif} \\
\arm{E3} -- \arm{E2} & fogalomgráf & RQ2 & -- \\
\arm{E2h} -- \arm{E2} & hibrid visszakeresés & RQ2 & -- \\
\arm{E2f} -- \arm{E2} & vak megoldási próba & RQ2 & -- \\
\arm{E1}, \arm{E2} -- \arm{B-doc-L} & módszer vs.\ status quo (helyi) & RQ1 & \ref{h:lefed} \\
\arm{E4b} -- \arm{B-doc-S} & módszer vs.\ status quo (szerver) & RQ1 & \ref{h:lefed} \\
\arm{E2} -- \arm{E4} & helyi módszer vs.\ szerver alapvonal & RQ3 & \ref{h:helyi} \\
\arm{E4} -- \arm{E0} & modellméret & RQ3 & \ref{h:helyi} \\
\arm{E5a/b/c} -- \arm{E0} & darabolási stratégia & RQ2 & -- \\
\bottomrule
\end{tabular}
\end{table}

\section{Mutatók}

A mutatók három csoportba sorolhatók: bírót nem igénylő \emph{automatikus} mutatók, a bíró
ítéletein alapuló \emph{tartalmi} mutatók, valamint \emph{költség}mutatók.

\subsection{Automatikus mutatók}

A formai megfelelés ($F$, \eqref{eq:format}), a lefedettség ($\mathrm{Cov}$, \eqref{eq:cov}), a
szivárgás és a nem-kérdés arány ($L$, $N$, \eqref{eq:leak}), valamint a duplikátumráta definícióját \az{\ref{ch:formalizmus}}.~fejezet adta meg. Ezeken felül mérjük a meta-hivatkozás arányát (a törzs
„a szöveg szerint”, „a fenti részlet” jellegű fordulatot tartalmaz-e, reguláris kifejezéssel), a kulcs
lexikális alátámasztottságát (a kulcs szótokenjeinek a kontextusban előforduló aránya), valamint két, a
referenciakérdésekre épülő mutatót:
\begin{align}
\mathrm{GoldSim}(\Exam)&=\frac{1}{|\Exam|}\sum_{q\in\Exam}\max_{g\in\mathcal{R}}\cos\bigl(\phi(q),\phi(g)\bigr),\\
\mathrm{GoldRecall}(\Exam)&=\frac{1}{|\mathcal{R}|}\sum_{g\in\mathcal{R}}\1\Bigl[\max_{q\in\Exam}\cos\bigl(\phi(q),\phi(g)\bigr)\ge0{,}85\Bigr],
\end{align}
ahol $\mathcal{R}$ a dokumentum referenciakérdéseinek halmaza. A visszakeresés minőségét egy
\emph{próba} méri: minden referenciakérdésre lefuttatjuk a visszakeresést, és megnézzük, hányadik
helyen szerepel az első olyan darab, amely a bizonyító mondatot tartalmazza ($\mathrm{rank}_g$).
Ebből
\begin{equation}
\mathrm{Recall@}k=\frac{1}{|\mathcal{R}|}\sum_{g}\1[\mathrm{rank}_g\le k],\qquad
\mathrm{MRR}=\frac{1}{|\mathcal{R}|}\sum_{g}\frac{1}{\mathrm{rank}_g},
\end{equation}
$k\in\{1,3,5,10\}$ mellett, ahol a meg nem talált bizonyítékhoz $1/\mathrm{rank}_g=0$ tartozik.

\subsection{Tartalmi mutatók}

A tartalmi mutatók a bíró (\ref{sec:biro}.~szakasz) ítéletein alapulnak: az alátámasztottság ($G$), a
disztraktor-validitás ($V$), a vak megoldhatóság ($A$) és ezekből a minőségi index ($Q$,
\eqref{eq:qindex}), valamint négy 1--5 skálájú rubrikapontszám és egy bináris „felhasználnám”-ítélet
(\ref{tab:rubrika}.~táblázat). A rubrika szó szerint azonos a bíró és az oktatók számára, ami a bíró
validálásának (\ref{h:meres}.~hipotézis) előfeltétele.

\begin{table}[htbp]
\centering
\caption{A bíró és az oktatók közös értékelési rubrikája (rövidítve; a teljes szöveg \az{\ref{app-rubrika}}.~függelékben).}
\label{tab:rubrika}
\small
\begin{tabularx}{\textwidth}{@{}>{\raggedright\arraybackslash}p{2.6cm}>{\raggedright\arraybackslash}X@{}}
\toprule
Szempont & Horgonyok \\
\midrule
Helyesség & 1 = hibás kulcs vagy két helyes opció; 3 = helyes, de pontatlan; 5 = teljesen helyes, egyértelmű \\
Érthetőség & 1 = zavaros; 3 = nehezen érthető; 5 = vizsgakész megfogalmazás, nincs „a szöveg szerint” \\
Disztraktorok & 1 = abszurd vagy szintén helyes; 3 = részben hihető; 5 = mind hihető és hibás \\
Szintillesztés & 1 = egyértelműen rossz Bloom-szint; 3 = határeset; 5 = illik a kért nehézséghez \\
Felhasználnám & igen/nem: bekerülne-e egy valódi dolgozatba legfeljebb apró javítással \\
\bottomrule
\end{tabularx}
\end{table}

\subsection{Költségmutatók}

Vizsgánként rögzítjük a teljes és a szakaszonkénti futási időt (kinyerés, darabolás, indexelés,
tervezés, generálás), a nyelvimodell-hívások számát szakaszonként, a prompt- és a választokenek számát,
az értelmezhetetlen válaszok és a hálózati hibák számát, a GPU-memória alapszintjét és csúcsát, valamint
azt, hogy a modell teljes egészében a GPU-n futott-e (az Ollama \code{/api/ps} végpontja alapján; ha nem,
a modell egy része a CPU-ra került, ami a mérés érvényességét érinti).

\section{A bíró}
\label{sec:biro}

A bíró a gpt-oss-120b~\cite{gptoss2025} modell az egyetemi GenAI-szerveren, $T=0$ hőmérséklettel,
alacsony gondolkodási szinttel („Reasoning: low”) és 16\,384 tokenes kontextussal. A választás
elsődleges szempontja, hogy más modellcsaládba tartozzon, mint bármelyik generátor (Qwen, Llama), így
az önpreferencia torzítása~\cite{panickssery2024selfpref} nem lép fel. Kérdésenként két hívást végez:

\begin{enumerate}
\item \textbf{Áttekintés} (látja a kulcsot): igazolja-e a forrás a megjelölt kulcsot, egy szó szerinti
bizonyító mondat idézésével; minden disztraktorra, hogy a forrás szerint hamis-e és témába vágó-e;
a négy rubrikapontszám; a kérdés Bloom-szintje; és a „felhasználnám”-ítélet, rövid indoklással.
\item \textbf{Vak megoldás} (nem látja a kulcsot): a forrás alapján válaszolnia kell a kérdésre,
véletlenül permutált opciókkal.
\end{enumerate}

A bíró forrásként alapértelmezésben a kérdéshez tartozó darabokat kapja, nem a teljes dokumentumot. Ezt
a fő vizsgálat futási idejének csökkentése indokolja (a bíró egy hívása a megosztott szerveren kb.\
14 másodperc, és ez a mérés szűk keresztmetszete); a választás korlátja, hogy a bíró a dokumentum más
részeiben lévő ellentmondást nem láthatja (\ref{sec:validitas}.~szakasz). A \code{context: document}
beállítás a teljes dokumentumot adja át, a pilot így futott.

\section{Oktatói értékelés}

A bíró csak akkor tekinthető érvényes mérőeszköznek, ha ítéletei az emberi szakértőkével egyeznek. A
fő vizsgálatban ezért három oktató vakon értékel egy rétegzett mintát: a négy fő karból
(\arm{B-doc-S}, \arm{E0}, \arm{E1}, \arm{E2} vagy \arm{E2x}) karonként 15 kérdést, azaz oktatónként 60
kérdést, valamint 10 közös kalibrációs kérdést. A lapok nem tartalmazzák a kar nevét, a kérdések
sorrendje véletlen, a rubrika azonos a bíróéval. Az értékelő--értékelő egyezést Krippendorff-$\alpha$-val
(ordinális skálán), a bíró--oktató egyezést szempontonként Spearman-$\rho$-val és Krippendorff-$\alpha$-val
mérjük, a bináris ítéletekre Cohen-$\kappa$-t számolunk.

\section{Statisztikai elemzés}
\label{sec:statisztika}

\subsection{Az elemzés egysége}

A kérdések nem függetlenek: ugyanabból a dokumentumból, ugyanazon a futáson belül készülnek, így a
kérdésszintű tesztek a szabadsági fokot mesterségesen felfújnák (pszeudoreplikáció). Ezért minden
mutatót először dokumentumon belül átlagolunk (a kérdésekre és a magokra), és a statisztikai
következtetés egysége a \emph{dokumentum}. Egy $m$ mutató $a$ karon mért értéke a $d$ dokumentumon
\begin{equation}
\bar m_{a,d}=\frac{1}{|\Omega_a|}\sum_{\omega\in\Omega_a}m\bigl(\Exam_a(d,\omega)\bigr),
\end{equation}
ahol $\Omega_a$ a kar magjainak halmaza.

\subsection{Konfidenciaintervallum}

A karonkénti átlag 95\,\%-os konfidenciaintervallumát percentilis bootstrappel~\cite{efron1993}
becsüljük, a dokumentumok visszatevéses újramintavételezésével, $B=10\,000$ ismétléssel:
\begin{equation}
\hat\mu^{*(b)}=\frac{1}{D}\sum_{i=1}^{D}\bar m_{a,d_i^{*(b)}},\qquad
\mathrm{CI}_{95}=\Bigl[\hat\mu^{*}_{(0{,}025)},\;\hat\mu^{*}_{(0{,}975)}\Bigr],
\end{equation}
ahol $d_i^{*(b)}$ a $b$-edik mintában a $D$ dokumentumból egyenletesen, visszatevéssel húzott elem.

\subsection{Páros összehasonlítás}

Két kar összevetésére a páros Wilcoxon-féle előjeles rangpróbát~\cite{wilcoxon1945} használjuk, amely
nem feltételez normalitást -- ami a $[0,1]$-be szorított, gyakran plafonhatást mutató rátáknál
lényeges. A $d$ dokumentumon mért különbség $\Delta_d=\bar m_{b,d}-\bar m_{a,d}$; a nullától különböző
különbségeket abszolút értékük szerint rangsoroljuk ($R_d$), és
\begin{equation}
W^{+}=\sum_{d:\Delta_d>0}R_d,\qquad W^{-}=\sum_{d:\Delta_d<0}R_d .
\end{equation}
A hatásméret a páros rangbiszeriális korreláció~\cite{kerby2014}:
\begin{equation}
r_{rb}=\frac{W^{+}-W^{-}}{W^{+}+W^{-}}\in[-1,1],
\end{equation}
amely az előjeles „egyszerű különbség” értelmezést adja: a kedvező és kedvezőtlen rangtömeg különbsége.

\begin{megjegyzes}[A minimális mintaméret]
Egzakt kétoldali Wilcoxon-próbánál a legkisebb elérhető $p$-érték $D$ nem nulla különbségű pár esetén
$2\cdot 2^{-D}$ (amikor minden különbség azonos előjelű). $D=5$-re ez $0{,}0625$, $D=6$-ra $0{,}03125$.
Hat dokumentumnál kevesebből tehát $p<0{,}05$ \emph{elvileg sem} érhető el; az egyetlen dokumentumos
pilot eredményei ezért kizárólag leíró jellegűek.
\end{megjegyzes}

\subsection{Többszörös összehasonlítás}

\Az{\ref{tab:osszevetesek}}.~táblázat összehasonlításainak $p$-értékeit Holm szekvenciális
eljárásával~\cite{holm1979} korrigáljuk. Ha a rendezett $p$-értékek $p_{(1)}\le\dots\le p_{(M)}$, a
korrigált értékek
\begin{equation}
\tilde p_{(i)}=\max_{j\le i}\;\min\bigl(1,\;(M-j+1)\,p_{(j)}\bigr),
\end{equation}
amely a családonkénti elsőfajú hibát $\alpha$ szinten tartja, a Bonferroni-korrekciónál nagyobb
erővel.

\subsection{Egyezési mutatók}

Két értékelő (vagy az ellenőrző és a bíró) bináris ítéleteinek egyezését Cohen-$\kappa$-val
mérjük~\cite{cohen1960kappa}:
\begin{equation}
\kappa=\frac{p_o-p_e}{1-p_e},\qquad p_e=\sum_{c}\hat\pi_{1,c}\,\hat\pi_{2,c},
\end{equation}
ahol $p_o$ a megfigyelt egyezési arány, $\hat\pi_{j,c}$ pedig a $j$-edik értékelő $c$ kategóriájának
relatív gyakorisága. Több értékelő és hiányzó értékek esetén Krippendorff-$\alpha$-t
használunk~\cite{krippendorff2004}:
\begin{equation}
\alpha=1-\frac{D_o}{D_e}=1-(n-1)\frac{\sum_{c,k}o_{ck}\,\delta^2_{ck}}{\sum_{c,k}n_c\,n_k\,\delta^2_{ck}},
\end{equation}
ahol $o_{ck}$ a koincidenciamátrix, $n_c$ a $c$ érték marginális gyakorisága, $n$ a párosítható
értékek száma, ordinális skálán pedig
$\delta^2_{ck}=\bigl(\sum_{g=c}^{k}n_g-\tfrac{n_c+n_k}{2}\bigr)^2$. Az irodalmi konvenció szerint
$\alpha\ge0{,}8$ megbízható, $0{,}667\le\alpha<0{,}8$ óvatos következtetésre elegendő.

\subsection{Az ellenőrző mint osztályozó}

Az ellenőrző lánc egy bináris osztályozó: minden kérdést \emph{átengedett} vagy \emph{megjelölt}
kategóriába sorol. A bíró alátámasztottsági ítéletét referenciának tekintve az ellenőrző
\begin{equation}
\mathrm{Prec}=\frac{|\{q:\text{átengedett}\wedge G(q)=1\}|}{|\{q:\text{átengedett}\}|},\qquad
\mathrm{Rec}=\frac{|\{q:\text{átengedett}\wedge G(q)=1\}|}{|\{q:G(q)=1\}|}
\end{equation}
pontossággal és felidézéssel, valamint a két bináris változó közötti Cohen-$\kappa$-val jellemezhető.
Egy hasznos ellenőrzőtől $\kappa\gg0$ várható; $\kappa\approx0$ azt jelenti, hogy az ellenőrző
ítélete a bíróétól független, vagyis a jelzése az oktató számára nem hordoz információt.

\section{Mérési validitás: a korábbi mérések hibái}
\label{sec:meresi-hibak}

A keretrendszer építése két olyan hibát tárt fel, amely a rendszer korábbi, végponttól végpontig tartó
méréseit érvénytelenítette. Ezek bemutatása azért lényeges, mert pontosan azt a kockázatot illusztrálják,
amely ellen a módszertan védekezik: egy mérés, amely „működik”, mégis mást mér, mint amit gondolunk.

\paragraph{Tartalékvizsgák eredményként.} Ha minden modellhívás sikertelen volt, a háttérrendszer egy
beégetett, egykérdéses tartalékvizsgát adott vissza hibaüzenet helyett. Egy 2026. május 3-i tesztfutás
mindhárom tesztdokumentumot 4--9 másodperc alatt „generáltnak” jelentette, holott mindhárom „vizsga”
ez a tartalék volt. A keretrendszer azóta a tartalékot hibajelzők alapján felismeri, és sikertelen
futásként számolja.

\paragraph{Csonkolt darabolás.} \Az{\ref{sec:csonkolt}}.~szakaszban elemzett hiba miatt a dokumentumok
első oldala utáni részét a darabolás gyakorlatilag a tartalomtól függetlenül vágta. A kódoló azóta
ablakozva dolgozik; az \arm{E5a} kar a régi viselkedést méri.

\paragraph{Válaszszivárgás.} A harmadik hibát a pilot tárta fel: a kérdés törzsébe másolt kulcs a
bíró és az automatikus mutatók értékét is felfelé torzítja (\ref{sec:szivargas}.~szakasz).

\section{Hardver}

A helyi karok egy fogyasztói GPU-n (NVIDIA GeForce RTX 3070 Ti, 8\,GB) és -- a hosszú futások miatt -- egy
egyetemi munkaállomáson (NVIDIA RTX A4500) futnak; mivel minden futás rögzíti a GPU típusát, a két gépen
készült eredmények elkülöníthetők, és a futásidők csak azonos hardveren hasonlíthatók össze. A fő
vizsgálat futtatási tervét (füstteszt, helyi és szerverkarok, pontozás, oktatói értékelés, jelentés) \az{\ref{sec:futasterv}}.~függelék tartalmazza.
